%==========================================================
% DOCUMENTCLASS
%==========================================================
\documentclass[12pt,oneside]{book}

%==========================================================
% METADATOS
%==========================================================
\newcommand{\universidad}{Universidad de Buenos Aires (UBA)}
\newcommand{\facultad}{Facultad de Derecho}
\newcommand{\materia}{Derecho Constitucional Profundizado}
\newcommand{\titulo}{Segunda parte: Derechos políticos y electorales}
\newcommand{\autor}{Isaac Edgar Camacho Ocampo}
\newcommand{\anio}{2026}
\newcommand{\reflexion}{Comprender cada concepto es el primer paso para convertir el estudio en conocimiento.}

%==========================================================
% PAQUETES
%==========================================================
\usepackage[utf8]{inputenc}
\usepackage[T1]{fontenc}
\usepackage[spanish,es-nodecimaldot]{babel}

\usepackage[
    top=2.5cm,
    bottom=2.5cm,
    left=3cm,
    right=2.5cm,
    headheight=15pt
]{geometry}

\usepackage{amsmath}
\usepackage{amssymb}
\usepackage{mathtools}

\usepackage{graphicx}
\usepackage{float}
\usepackage{caption}
\usepackage{subcaption}

\usepackage{booktabs}
\usepackage{array}
\usepackage{longtable}
\usepackage{tabularx}

\usepackage{enumitem}

\usepackage{xcolor}
\usepackage[most]{tcolorbox}

\usepackage{fancyhdr}
\usepackage{ragged2e}

\usepackage[
    colorlinks=true,
    linkcolor=blue!50!black,
    citecolor=green!40!black,
    urlcolor=blue!60!black,
    pdftitle={\titulo},
    pdfauthor={\autor},
    pdfsubject={Apuntes universitarios},
    bookmarks=true
]{hyperref}

%==========================================================
% CONFIGURACIÓN
%==========================================================
\graphicspath{
    {./img/}
    {./graficos/}
    {./figuras/}
}

\setcounter{tocdepth}{2}

\setlength{\parindent}{0pt}
\setlength{\parskip}{0.5em}

\setlist[itemize]{
    topsep=0.3em,
    itemsep=0.2em
}

\setlist[enumerate]{
    topsep=0.3em,
    itemsep=0.2em
}

\clubpenalty=10000
\widowpenalty=10000

\pagestyle{fancy}
\fancyhf{}
\fancyhead[L]{\nouppercase{\leftmark}}
\fancyhead[R]{\materia}
\fancyfoot[C]{\thepage}
\renewcommand{\headrulewidth}{0.4pt}
\renewcommand{\footrulewidth}{0pt}

\fancypagestyle{plain}{
    \fancyhf{}
    \fancyfoot[C]{\thepage}
    \renewcommand{\headrulewidth}{0pt}
}

%==========================================================
% COMANDOS
%==========================================================
\newcommand{\abs}[1]{\left|#1\right|}
\newcommand{\norm}[1]{\left\lVert#1\right\rVert}
\newcommand{\vect}[1]{\mathbf{#1}}
\newcommand{\deriv}[2]{\frac{d#1}{d#2}}
\newcommand{\pderiv}[2]{\frac{\partial#1}{\partial#2}}

%==========================================================
% ENTORNOS / CAJAS
%==========================================================
\newtcolorbox{definition}[1]{
    enhanced,
    breakable,
    title={Definición: #1},
    fonttitle=\bfseries,
    colback=blue!3,
    colframe=blue!50!black,
    boxrule=0.8pt,
    arc=2mm
}

\newtcolorbox{important}{
    enhanced,
    breakable,
    title={Importante},
    fonttitle=\bfseries,
    colback=yellow!8,
    colframe=orange!70!black,
    boxrule=0.8pt,
    arc=2mm
}

\newtcolorbox{example}[1]{
    enhanced,
    breakable,
    title={Ejemplo: #1},
    fonttitle=\bfseries,
    colback=green!4,
    colframe=green!40!black,
    boxrule=0.8pt,
    arc=2mm
}

\newtcolorbox{formula}[1]{
    enhanced,
    breakable,
    title={Fórmula / Regla: #1},
    fonttitle=\bfseries,
    colback=purple!4,
    colframe=purple!50!black,
    boxrule=0.8pt,
    arc=2mm
}

\newtcolorbox{exercise}[1]{
    enhanced,
    breakable,
    title={Ejercicio: #1},
    fonttitle=\bfseries,
    colback=gray!5,
    colframe=gray!60!black,
    boxrule=0.8pt,
    arc=2mm
}

\newtcolorbox{note}{
    enhanced,
    breakable,
    title={Nota},
    fonttitle=\bfseries,
    colback=white,
    colframe=gray!50,
    boxrule=0.6pt,
    arc=2mm
}

%==========================================================
% CONFIGURACIÓN FINAL
%==========================================================
\renewcommand{\arraystretch}{1.25}

%==========================================================
% DOCUMENT
%==========================================================
\begin{document}

%----------------------------------------------------------
% PORTADA
%----------------------------------------------------------
\begin{titlepage}

\thispagestyle{empty}

\begin{center}

\vspace*{1cm}

{\large \textsc{\universidad}}

\vspace{0.2cm}

{\large \textsc{\facultad}}

\vspace{3cm}

{\Huge \textbf{\materia}}

\vspace{1cm}

{\LARGE \titulo}

\vspace{0.8cm}

\textit{\reflexion}

\vspace{1.2cm}

\rule{0.70\textwidth}{0.5pt}

\vspace{0.5cm}

{\large Apuntes de estudio}

\vspace{0.5cm}

\rule{0.70\textwidth}{0.5pt}

\vfill

{\large \autor}

\vspace{0.3cm}

{\large \anio}

\end{center}

\vfill

\begin{flushleft}
\textit{Este es un modesto aporte a los estudiantes de ``\materia'' no pretende sustituir el material oficial de la materia.}
\end{flushleft}

\end{titlepage}

%----------------------------------------------------------
% ÍNDICE
%----------------------------------------------------------
\frontmatter
\tableofcontents
\mainmatter

%----------------------------------------------------------
% INTRODUCCIÓN
%----------------------------------------------------------
\chapter*{Introducción: alcance y relevancia del tema}
\addcontentsline{toc}{chapter}{Introducción: alcance y relevancia del tema}
\markboth{Introducción}{Introducción}

\section*{Temas tratados}

Estos apuntes cierran la segunda parte de la materia con los \textbf{derechos políticos y electorales}. Se recorre el diseño constitucional del sistema político (república, representación, elección directa e indirecta de autoridades); la reserva legislativa en materia electoral y de partidos, con la prohibición de los decretos de necesidad y urgencia; el rol de los partidos políticos y el sufragio, a través de los fallos \textbf{Ríos} y \textbf{Mignone}; y el principio de \textbf{alternancia}, a partir del fallo de la Corte sobre la reelección indefinida en Formosa. El desarrollo concluye con dos prácticas cuestionadas: la \textbf{ley de lemas} y las \textbf{candidaturas testimoniales}.

\section*{Relevancia profesional del tema}

\begin{itemize}
    \item \textbf{Litigio y asesoramiento electoral:} permite saber quién puede nominar candidatos, qué exige la ley a un partido para existir y cómo se discute una candidatura.
    \item \textbf{Argumentación sólida:} Ríos, Mignone, Formosa y la OC-28/21 son precedentes que se citan en escritos, dictámenes y exámenes.
    \item \textbf{Control de constitucionalidad:} permite medir una norma, por ejemplo una cláusula provincial de reelección, contra los principios republicano y representativo.
    \item \textbf{Trabajo en el Estado o el Congreso:} explica por qué ciertas materias solo se regulan por ley y no por decreto de necesidad y urgencia, clave para evaluar la validez de una norma.
    \item \textbf{Razonamiento por analogía:} el debate sobre los intendentes enseña a extender un precedente a un caso nuevo, una habilidad central de todo abogado.
\end{itemize}

\section*{Utilidad más allá del ejercicio profesional}

El conocimiento de estos temas resulta útil incluso para quien no ejerce el Derecho. En Argentina el voto es obligatorio: saber qué se elige, quién puede postularse y con qué límites permite votar con criterio. Los derechos políticos no desaparecen por estar detenido sin condena firme, y conocerlo protege a cualquiera. La alternancia y la división de poderes impiden que alguien concentre poder indefinidamente, algo que impacta en libertades, economía y servicios cotidianos. Además, el tema entrena el pensamiento crítico para distinguir una regla legítima de un atajo (un decreto de necesidad y urgencia electoral, una lista testimonial, un lema) que confunde al votante.

\section*{Relación entre los temas: un ejemplo cotidiano}

La elección de la comisión directiva de un club de barrio permite ilustrar los cuatro bloques temáticos como cuatro preguntas que se plantearían en esa elección.

\begin{table}[H]
\centering
\begin{tabularx}{\textwidth}{>{\raggedright\arraybackslash}X >{\raggedright\arraybackslash}X}
\toprule
\textbf{En un club de barrio\ldots} & \textbf{En estos apuntes\ldots} \\
\midrule
El estatuto fija quién elige a quién, y solo la asamblea puede modificarlo. &
Capítulo 1: el Ejecutivo no puede cambiar las reglas electorales por decreto; solo el Congreso, mediante ley. \\
\addlinespace
Solo pueden presentar lista las agrupaciones reconocidas, y todos los socios sin sanción firme conservan su voto. &
Capítulo 2: los partidos son el único vehículo de nominación (Ríos) y el detenido sin condena firme sigue votando (Mignone). \\
\addlinespace
El presidente del club no puede quedarse 30 años en el cargo. &
Capítulo 3: la alternancia y el fallo Formosa contra la reelección indefinida. \\
\addlinespace
Las trampas: listas que suman votos entre sí y candidatos que figuran sabiendo que no asumirán. &
Capítulo 4: la ley de lemas y las candidaturas testimoniales. \\
\bottomrule
\end{tabularx}
\end{table}

%----------------------------------------------------------
% CAPÍTULO 1
%----------------------------------------------------------
\chapter[Diseño constitucional del sistema político]{Diseño constitucional del sistema político}

\section{Introducción a los derechos políticos y electorales}

Los derechos políticos y electorales constituyen el último tema novedoso que se introduce en esta segunda parte de la materia. Se trata, como tantas otras cuestiones del derecho constitucional, de un área que es una materia en sí misma y un gran campo del derecho, por lo que se abordan únicamente temas muy concretos.

El propósito es ofrecer un panorama más profundizado que el de la parte de elementos, articulando los \textbf{principios constitucionales más importantes} y analizando los fallos en materia electoral y de representación política más significativos de la historia moderna de la Corte, es decir, de los últimos 25 o 30 años.

\section{El sistema republicano y la elección de autoridades}

\subsection{Carácter republicano y elección directa}

En materia de representación política y electoral, la organización constitucional argentina responde a un sistema \textbf{republicano}. Existe, por lo tanto, una serie de cargos públicos para los cuales los ciudadanos pueden emitir el sufragio y elegirlos. La \textbf{elección directa} corresponde a la alta jerarquía del Poder Ejecutivo y a los miembros del Congreso de la Nación.

\subsection{Elección indirecta de los magistrados}

Los magistrados se designan de forma \textbf{indirecta}, mediante un sistema que la actual Constitución estableció a través de la intermediación del \textbf{Consejo de la Magistratura}. Esto alcanza a los magistrados federales y a los miembros de la Corte.

Asimismo, el Congreso de la Nación, mediante la elevación de ternas y el posterior acuerdo del Senado, aún hoy designa los jueces nacionales con asiento en la Capital Federal. Sin embargo, de acuerdo con las últimas iniciativas de la gestión política actual y con el acuerdo de transferencia, esa justicia nacional se modificará paulatinamente hacia una justicia eminentemente local, particularmente en el fuero laboral. En consecuencia, intervendrán en la designación de los futuros jueces nacionales (en lo comercial, civil y del trabajo) la Legislatura local y el Consejo de la Magistratura local.

\begin{note}
\textbf{Nota editorial de los apuntes originales.} Con precisión técnica: los jueces de la Corte Suprema son designados por el Poder Ejecutivo con acuerdo del Senado (art.~99, inc.~4, CN); los jueces federales inferiores se seleccionan mediante el Consejo de la Magistratura, que eleva ternas vinculantes al Poder Ejecutivo, con acuerdo del Senado (arts.~99, inc.~4, y 114, CN).
\end{note}

\section{Reserva legislativa en materia electoral y prohibición de decretos de necesidad y urgencia}

\subsection{El constituyente y la regulación electoral}

Se advierte un especial celo del constituyente en dotar de un carácter eminentemente legislativo a la materia electoral y de partidos políticos: la regulación del sistema de partidos y de la materia electoral se asienta exclusivamente en el Congreso de la Nación. La razón es que, de no ser así, cualquier poder de turno podría modificarla fácilmente para su conveniencia.

Como consecuencia de esta necesidad de que solo el Congreso la regule, se requiere una \textbf{ley}. Es una de las \textbf{cuatro materias vedadas} para los decretos de necesidad y urgencia: en materia electoral y de partidos políticos, el Poder Ejecutivo no puede emitir decretos de necesidad y urgencia.

\begin{formula}{Art.~99, inc.~3, tercer párrafo --- Constitución Nacional}
«El Poder Ejecutivo no podrá en ningún caso bajo pena de nulidad absoluta e insanable, emitir disposiciones de carácter legislativo. Solamente cuando circunstancias excepcionales hicieran imposible seguir los trámites ordinarios previstos por esta Constitución para la sanción de las leyes, y no se trate de normas que regulen materia penal, tributaria, electoral o el régimen de los partidos políticos, podrá dictar decretos por razones de necesidad y urgencia, los que serán decididos en acuerdo general de ministros que deberán refrendarlo, conjuntamente con el jefe de gabinete de ministros.»
\end{formula}

\subsection{Ejemplo: la ley de primarias}

\begin{example}{modificación de las primarias simultáneas y obligatorias}
Si bien existen intenciones concretas de modificar la ley que establece las primarias simultáneas y obligatorias, el Poder Ejecutivo no podría hacerlo nunca por decreto. Esa vía sería cuestionable aun sin la prohibición expresa, pero la prohibición disipa cualquier duda: el Poder Ejecutivo no puede dictar regulaciones mediante decretos de necesidad y urgencia en materia electoral y de partidos políticos.
\end{example}

\begin{note}
\textbf{Nota editorial de los apuntes originales.} La ley de primarias abiertas, simultáneas y obligatorias (PASO) es la Ley 26.571 (2009). Al ser materia electoral, solo puede ser modificada por ley del Congreso.
\end{note}

\begin{important}
\textbf{Reserva de ley electoral.} La regulación electoral y de partidos políticos es exclusiva del Congreso para impedir que el poder de turno la modifique a su conveniencia; por eso el Poder Ejecutivo no puede legislar sobre ella mediante decretos de necesidad y urgencia.
\end{important}

%----------------------------------------------------------
% CAPÍTULO 2
%----------------------------------------------------------
\chapter[Partidos políticos y sufragio]{Partidos políticos y sufragio}

\section{Los partidos políticos como instituciones del sistema democrático}

\subsection{Representación y rol de los partidos}

El sistema argentino es también \textbf{representativo y federal}, inserto en la división de poderes del régimen republicano. El pueblo no delibera ni gobierna sino a través de sus representantes, y los partidos políticos son instituciones clave en el desarrollo del sistema político de la Nación. Por eso existe la \textbf{Ley Orgánica de Partidos Políticos}, que regula las condiciones en las cuales estas instituciones pueden funcionar, y a ellas les está concedida \textbf{en exclusividad} la posibilidad de nominar hombres y mujeres a cargos públicos electivos.

\begin{formula}{Art.~22 --- Constitución Nacional}
«El pueblo no delibera ni gobierna, sino por medio de sus representantes y autoridades creadas por esta Constitución. Toda fuerza armada o reunión de personas que se atribuya los derechos del pueblo y peticione a nombre de éste, comete delito de sedición.»
\end{formula}

\begin{formula}{Art.~38 --- Constitución Nacional}
«Los partidos políticos son instituciones fundamentales del sistema democrático. Su creación y el ejercicio de sus actividades son libres dentro del respeto a esta Constitución, la que garantiza su organización y funcionamiento democráticos, la representación de las minorías, la competencia para la postulación de candidatos a cargos públicos electivos, el acceso a la información pública y la difusión de sus ideas. El Estado contribuye al sostenimiento económico de su actividad y de la capacitación de sus dirigentes. Los partidos políticos deberán dar publicidad del origen y destino de sus fondos y patrimonio.»
\end{formula}

\subsection{Requisitos de la Ley Orgánica de Partidos Políticos}

Los partidos políticos deben obtener una \textbf{personería jurídica} y organizarse mediante un \textbf{estatuto}, y están controlados por la justicia electoral. La justicia electoral en el ámbito federal regula el funcionamiento de los partidos y ejerce el control de las condiciones de aptitud para presentarse a elecciones con sus candidatos nominados. La Ley Orgánica establece cuatro o cinco principios vectores.

En primer lugar, el partido debe ser una persona jurídica reconocida en el campo jurídico. En segundo lugar, debe contar con un número de miembros, los afiliados: para ser reconocido como tal debe tener cuanto menos un piso del 4\,\% del padrón electoral en cada jurisdicción. A nivel nacional existen partidos con representación por jurisdicción, es decir, por cada una de las provincias y la Ciudad de Buenos Aires; en cada jurisdicción, si cuentan con ese piso de afiliados, pueden acceder al reconocimiento de la persona jurídica y a la presentación en elecciones. Como el piso es por jurisdicción, puede haber partidos con reconocimiento solo en algunas de ellas. Los partidos con más de cinco jurisdicciones reconocidas tienen estatus de partido a nivel nacional.

En tercer lugar, el partido debe tener una organización interna mediante un estatuto y un régimen de renovación y alternancia de sus autoridades. Al tratarse de organizaciones que propenden al sistema representativo y democrático, deben contar internamente con un modo de organización que asegure la periodicidad en la elección de sus cargos directivos.

\begin{table}[H]
\centering
\begin{tabularx}{\textwidth}{>{\raggedright\arraybackslash}p{0.24\textwidth} >{\raggedright\arraybackslash}X}
\toprule
\textbf{Requisito} & \textbf{Contenido} \\
\midrule
\textbf{Personería jurídica} & Debe ser una persona jurídica reconocida en el campo jurídico. \\
\textbf{Afiliados} & Número mínimo de miembros afiliados («un piso») por jurisdicción para acceder al reconocimiento y a la presentación en elecciones. \\
\textbf{Organización interna} & Estatuto y régimen de renovación y alternancia de autoridades, que asegure la periodicidad en la elección de los cargos directivos. \\
\textbf{Alcance territorial} & El piso se exige por jurisdicción (provincias y Ciudad de Buenos Aires). Con más de cinco jurisdicciones reconocidas, hay estatus de partido nacional. \\
\textbf{Control} & La justicia electoral federal regula el funcionamiento y controla la aptitud para presentarse a elecciones con candidatos nominados. \\
\bottomrule
\end{tabularx}
\end{table}

% TODO: contenido incompleto o ambiguo en las notas originales (el porcentaje de afiliados se expone como 4 %; la nota editorial original lo corrige).
\begin{note}
\textbf{Nota editorial de los apuntes originales.} Corrección técnica: la Ley 23.298 exige la afiliación de un número de electores no inferior al cuatro por mil (4\,\textperthousand) del padrón del distrito (art.~7 bis), y no del 4\,\%. Para obtener el carácter de partido nacional se requiere el reconocimiento en cinco o más distritos (art.~8).
\end{note}

\section{El fallo Ríos: el monopolio partidario de la nominación}

El fallo Ríos tiene más de tres décadas y se inserta en la primera ola de fallos de derecho constitucional posteriores a la recuperación del sistema democrático en 1983, considerados fallos fundacionales. Consagra la legitimidad de los partidos políticos como el \textbf{único vehículo institucional} por el cual pueden nominarse los electores para los distintos cargos.

El razonamiento, de sentido común, es el siguiente: si se admitiera que cualquier sujeto fuera candidato por su sola voluntad, el sistema sería enormemente fragmentado y atomizado, muy probablemente inoperante para ejercer la legítima representación de la voluntad del pueblo. Por ello deben existir instituciones clave para el régimen electoral: los partidos políticos.

Frente a la pretensión de un particular de nominarse como candidato, la Corte sostuvo que no puede hacerlo si no es por vehículo de alguna de estas instituciones: \textbf{no puede haber una nominación independiente} a cargos públicos electivos.

\begin{note}
\textbf{Nota editorial de los apuntes originales.} «Ríos, Antonio Jesús» (CSJN, 1987). En el régimen vigente de entonces, la Corte sostuvo que la postulación de candidatos a cargos electivos nacionales corresponde a los partidos políticos.
\end{note}

\begin{important}
\textbf{Monopolio de la nominación.} Sin partidos como único vehículo de nominación, el sistema se fragmentaría y sería inoperante para representar la voluntad popular.
\end{important}

\section{Periodicidad electoral y el fallo Mignone}

\subsection{Periodicidad del sufragio}

Otro punto clave del sistema electoral es su periodicidad y habitualidad. Aun cuando suele ser objeto de críticas, el país tiene, a nivel general, \textbf{elecciones cada dos años}, a fin de asegurar la alternancia de los partidos políticos y el continuo ejercicio del sufragio, que es la fuente de la legitimidad de los poderes del Estado.

\subsection{El caso Mignone (2002): voto de las personas detenidas sin condena}

En el caso Mignone (2002), la Corte se pronunció sobre la regulación entonces existente, que \textbf{negaba el derecho al voto} a las personas detenidas sin condena firme, es decir, a quienes se encontraban con prisión preventiva. La Corte declaró la invalidez de esa disposición, y la ley posterior recepta este principio.

Hoy se entiende que privar del derecho al voto a quien todavía no tiene condena implica una violación del \textbf{principio de inocencia}. Esto es aún más relevante en el caso de personas privadas de su libertad, cuya capacidad para incidir sobre temas políticos debe ser extremadamente protegida, más allá de la inconducta que hayan podido cometer. La Corte afirmó, entonces, que los procesados no condenados deben tener asegurado su derecho al voto.

En la legislación actual, solo las condenas efectivas, firmes, por delitos dolosos y privativas de la libertad llevan como sanción adicional la imposibilidad de ejercer el derecho al voto.

\begin{table}[H]
\centering
\begin{tabularx}{\textwidth}{>{\raggedright\arraybackslash}p{0.38\textwidth} >{\raggedright\arraybackslash}X}
\toprule
\textbf{Situación de la persona} & \textbf{Derecho al voto} \\
\midrule
Procesada, detenida sin condena (prisión preventiva) & Debe tener asegurado el derecho al voto (fallo Mignone); privarla de él viola el principio de inocencia. \\
Con condena efectiva, firme, por delito doloso privativo de la libertad & Sanción adicional: imposibilidad de ejercer el derecho al voto (legislación actual). \\
\bottomrule
\end{tabularx}
\end{table}

\begin{important}
\textbf{Mignone.} El voto no puede quitarse a quien no tiene condena firme: hacerlo vulnera el principio de inocencia.
\end{important}

\section{Caracteres del sufragio}

El derecho al voto es \textbf{universal}, \textbf{secreto} y, en Argentina, \textbf{obligatorio}, lo que lo diferencia del de muchos otros países occidentales. El carácter tradicionalmente obligatorio del sufragio ha sido sostenido siempre, y es quizás uno de los puntos en que los partidos políticos no tienen disensos. El sufragio es, además, la \textbf{fuente de legitimación} de la designación de los cargos electivos.

\begin{formula}{Art.~37, primer párrafo --- Constitución Nacional}
«Esta Constitución garantiza el pleno ejercicio de los derechos políticos, con arreglo al principio de la soberanía popular y de las leyes que se dicten en consecuencia. El sufragio es universal, igual, secreto y obligatorio.»
\end{formula}

%----------------------------------------------------------
% CAPÍTULO 3
%----------------------------------------------------------
\chapter[Alternancia y reelección]{Alternancia y reelección}

\section{La alternancia en el poder}

La jurisprudencia ha destacado siempre la \textbf{alternancia} como un signo de afianzamiento de los sistemas democráticos, representativos y republicanos. La posibilidad de alternancia alimenta el sistema de frenos y contrapesos entre los distintos poderes del Estado.

\section{El fallo Formosa: reelección indefinida y artículo 132}

\subsection{Planteo del caso}

En un fallo reciente contra la provincia de Formosa, la Corte declaró la inconstitucionalidad sobreviniente del artículo 132 de la Constitución de esa provincia. El caso estaba dirigido a la situación del gobernador Gildo Insfrán. La Corte declaró la invalidez del artículo que disponía la reelección sin restricciones, es decir, indefinida, del gobernador y del vicegobernador.

\subsection{Fundamentos: la Opinión Consultiva 28/21}

La Corte recepta el criterio de la Corte Interamericana de Derechos Humanos expresado en la \textbf{Opinión Consultiva 28 de 2021}: la conservación indefinida del cargo ejecutivo desnaturaliza la estructura republicana que consagra la Constitución. Mantener el artículo de la Constitución provincial contraría, por tanto, el estándar que la Argentina sostiene.

Desde el punto de vista legal, la Corte hace mérito de los lineamientos y principios vectores del sistema político, pero se alimenta esencialmente de esa opinión consultiva. La Corte Interamericana tiene dos grandes competencias: la \textbf{competencia jurisdiccional}, en la que revisa eventuales violaciones al pacto, y las \textbf{opiniones consultivas}, que son decisiones interpretativas sobre el pacto. La Corte ha manifestado que las opiniones consultivas son también una fuente de hermenéutica del pacto, que forma parte de nuestra Constitución. Fue, por lo tanto, una cita legal de mucho peso, junto con el pacto de derechos políticos.

% TODO: contenido incompleto o ambiguo en las notas originales (el docente se refiere al «Pacto Internacional de Derechos Políticos»; la nota editorial original lo precisa).
\begin{note}
\textbf{Nota editorial de los apuntes originales.} Precisiones: (a) la OC-28/21 de la Corte IDH (7 de junio de 2021) versa sobre «La figura de la reelección presidencial indefinida en sistemas presidenciales»; (b) la competencia contenciosa de la Corte IDH no es propiamente una instancia de apelación, sino la resolución de casos por presuntas violaciones de la Convención Americana; (c) el instrumento citado es el Pacto Internacional de Derechos Civiles y Políticos, junto con la Convención Americana sobre Derechos Humanos; ambos tienen jerarquía constitucional (art.~75, inc.~22, CN).
\end{note}

\begin{formula}{Art.~23, inc.~1 --- Convención Americana sobre Derechos Humanos}
«Todos los ciudadanos deben gozar de los siguientes derechos y oportunidades: a) de participar en la dirección de los asuntos públicos, directamente o por medio de representantes libremente elegidos; b) de votar y ser elegidos en elecciones periódicas auténticas, realizadas por sufragio universal e igual y por voto secreto que garantice la libre expresión de la voluntad de los electores, y c) de tener acceso, en condiciones generales de igualdad, a las funciones públicas de su país.»
\end{formula}

\begin{formula}{Art.~25 --- Pacto Internacional de Derechos Civiles y Políticos}
«Todos los ciudadanos gozarán, sin ninguna de las distinciones mencionadas en el artículo 2, y sin restricciones indebidas, de los siguientes derechos y oportunidades: a) Participar en la dirección de los asuntos públicos, directamente o por medio de representantes libremente elegidos; b) Votar y ser elegidos en elecciones periódicas auténticas, realizadas por sufragio universal e igual y por voto secreto que garantice la libre expresión de la voluntad de los electores; c) Tener acceso, en condiciones generales de igualdad, a las funciones públicas de su país.»
\end{formula}

\begin{formula}{Art.~5 --- Constitución Nacional}
«Cada provincia dictará para sí una Constitución bajo el sistema representativo republicano, de acuerdo con los principios, declaraciones y garantías de la Constitución Nacional; y que asegure su administración de justicia, su régimen municipal, y la educación primaria. Bajo de estas condiciones el Gobierno federal, garante a cada provincia el goce y ejercicio de sus instituciones.»
\end{formula}

\begin{formula}{Art.~90 --- Constitución Nacional (modelo de reelección limitada)}
«El Presidente y vicepresidente duran en sus funciones el término de cuatro años y podrán ser reelegidos o sucederse recíprocamente por un solo período consecutivo. Si han sido reelectos o se han sucedido recíprocamente, no pueden ser elegidos para ninguno de ambos cargos, sino con el intervalo de un período.»
\end{formula}

\subsection{Alcance y fundamentos del fallo}

La Corte emite un fallo muy a la medida de la situación vigente. No sienta un estándar objetivo sobre cuántos períodos de reelección son admisibles. Sí deja aclarado que hasta dos reelecciones podrían ser toleradas, pero no emite una opinión concluyente sobre cuántos períodos generan la ilegitimidad.

Lo que sostiene es que existe una situación fáctica en la cual el gobernador llevaba más de 30 años en el poder (37, según la propia Corte) mediante reelecciones que, si bien eran a priori y formalmente legítimas, socavaron totalmente el sistema republicano. Ello se debe a que, mediante la continuidad en los cargos y la configuración de una suerte de \textbf{superestructura política permanente}, se debilita la alternancia.

La Corte afirma así que la situación, tal como se presenta fácticamente en la provincia, es reprochable desde el punto de vista constitucional, porque no es tolerable un período ininterrumpido semejante en el ejercicio del Ejecutivo provincial. Sin dar una precisión taxativa de cuál sería un período razonable, considera que lo observado es claramente opuesto a un régimen democrático razonable y, por lo tanto, declara la inconstitucionalidad del artículo de la Constitución provincial.

\begin{note}
\textbf{Nota editorial de los apuntes originales.} Fallo: «Confederación Frente Amplio Formoseño c/ Formosa, Provincia de s/ amparo» (CSJN, 19 de diciembre de 2024), que declaró la inconstitucionalidad del art.~132 de la Constitución de Formosa por contradecir el sistema republicano de los arts.~5, 123 y concordantes de la Constitución Nacional. La cifra de 37 años computa el tiempo total en la fórmula del Poder Ejecutivo provincial: Insfrán gobierna desde 1995 y antes fue vicegobernador.
\end{note}

\subsection{La reforma posterior del artículo 132}

La cuestión es hoy nuevamente objeto de estudio, porque la provincia modificó el artículo 132 y redactó uno muy parecido al de la Constitución Nacional. Sin embargo, insistió en que la actual gestión del gobernador debe considerarse la primera a los efectos de los dos períodos consecutivos que el artículo considera válidos. En los hechos, esto significa que Insfrán podría ser candidato por un período más. Es probable que en breve haya una aclaración o un fallo complementario.

\begin{important}
\textbf{Doctrina del fallo Formosa.} La reelección indefinida, aunque formalmente legítima, socava el sistema republicano y la alternancia. La Corte no fija un número exacto de períodos tolerables, pero toleraría hasta dos reelecciones y descalificó un período ininterrumpido como el del caso.
\end{important}

\section{Intendentes y otras provincias}

\subsection{El caso de los intendentes}

Un ejemplo de la discusión es el de los intendentes, cuya reelección era indefinida. Es una disputa política y una cuestión en debate público en el ámbito de la provincia: si se remueve nuevamente el artículo limitativo o se mantiene la idea de no más de una reelección. Aun cuando la intendencia pueda reputarse un cargo de nivel inferior (ciertamente, el más inferior con reconocimiento constitucional), plantea situaciones muy complejas, pues muchos intendentes han tenido períodos muy prolongados.

Habría que ver, en un fallo al efecto, si la Corte y los tribunales inferiores mantienen las restricciones al nivel de la doctrina de Formosa o se vuelven más laxos, entendiendo que un intendente no tiene un poder ni una entidad de derecho público tan contundente y que podría caber más laxitud.

\subsection{Perpetuación en las intendencias}

Aunque la intendencia sea el escalafón más bajo, hay partidos relativamente muy grandes manejados por la misma persona durante muchísimos años. En cualquier nivel de gobierno no es deseable una perpetuación de este tipo. Aunque no se trata de la misma entidad de cargos, el principio debería prevalecer.

% TODO: contenido incompleto o ambiguo en las notas originales (referencia a casos de Río Negro, frase parcialmente inteligible).

El fallo de Formosa sienta un precedente ordenador. La especulación política consiste en si se traslada de forma directa o no a los intendentes. Muy probablemente se necesite un nuevo fallo que analice la situación de los municipios. El concepto de fortalecimiento institucional por la alternancia del poder debería valer para todos los sistemas ejecutivos.

\subsection{El precedente Zamora (2013)}

En 2013 la Corte resolvió el caso de Gerardo Zamora, quien pretendía un nuevo mandato como gobernador de Santiago del Estero. Zamora había cumplido dos mandatos (2005--2009 y 2009--2013) y pretendía ser reelecto para un tercero, bajo el argumento de que el último período había sido elegido cuando se encontraba en modificación la nueva Constitución de Santiago del Estero. La Corte lo rechazó.

Es posible que ocurra algo similar con Insfrán, porque la Corte hizo mérito de la enorme extensión temporal. Aun cuando pudiera sostenerse que hay una legitimidad y que debe existir un período de adaptación (entendido como que la norma no se aplica de forma directa), esa deferencia resulta irrazonable cuando se trata de quien detenta el cargo desde hace más de 30 años. Si el período hubiera sido más razonable, la Corte quizás admitiría que ese período de adaptación pudiera tolerarse. Por lo tanto, este precedente indica que lo más probable es que haya un criterio igual. En el caso Zamora estaba también en discusión si se computaba o no el período en que había sucedido a la mujer, en cuanto a la inversión de los cargos de gobernador y vicegobernador.

\subsection{Provincias y regímenes de reelección}

Según lo señalado en el aula, Santa Cruz y Catamarca no tienen límite de reelección; Formosa era otra provincia en esa situación. San Juan prevé tres mandatos, pero con un límite. Respecto de Tucumán, el docente recordaba que había modificado el régimen con una reelección. En el caso Manzur, se señaló que se postuló como vicegobernador tucumano y, si resultaba electo, iniciaría su quinto período consecutivo como miembro de la fórmula gobernante; la Corte lo suspendió de forma cautelar y Manzur declinó su candidatura (2023).

% TODO: contenido incompleto o ambiguo en las notas originales (el docente y los alumnos no coinciden con precisión sobre el listado de provincias sin límite de reelección).
\begin{table}[H]
\centering
\begin{tabularx}{\textwidth}{>{\raggedright\arraybackslash}p{0.27\textwidth} >{\raggedright\arraybackslash}X}
\toprule
\textbf{Provincia} & \textbf{Lo expuesto} \\
\midrule
\textbf{Formosa} & Reelección indefinida (art.~132); declarada inconstitucional por la Corte. Reforma posterior: debate sobre si el período actual cuenta como el primero. \\
\textbf{Santa Cruz y Catamarca} & Régimen sin límite de reelección, según lo señalado. \\
\textbf{Tucumán} & Caso Manzur (2023): candidatura a vicegobernador, quinto período consecutivo en la fórmula gobernante; la Corte lo suspendió de forma cautelar y Manzur declinó. \\
\textbf{San Juan} & Prevé tres mandatos, pero con un límite. \\
\textbf{Santiago del Estero} & Caso Zamora (2013): la Corte rechazó la pretensión de un tercer mandato consecutivo. \\
\bottomrule
\end{tabularx}
\end{table}

\section{Cargos legislativos y comparación con Estados Unidos}

\subsection{¿Rige la alternancia para los legisladores?}

La alternancia no se aplica a los cargos legislativos. Quizás en la Argentina hubo algunos casos, pero no es algo muy pronunciado. En cambio, en Estados Unidos se discute desde hace tiempo la \textbf{perpetuación de los senadores}, que suelen permanecer en el cargo períodos muy largos, de 20 o 30 años.

El sistema estadounidense es bicameral como el argentino, aunque la forma de elección difiere bastante. Muchos estados tienen sistema uninominal y generalmente no se utiliza el sistema D’Hondt; hay una menor proporcionalidad y un mayor afianzamiento del bipartidismo y de las mayorías.

En la Argentina los diputados duran 4 años y los senadores 6; en Estados Unidos la diferencia entre las cámaras es más notoria, porque los representantes duran 2 años. Un mandato de senador equivale a tres de diputado en el caso estadounidense, y la disparidad entre las dos cámaras es mayor que en la Argentina.

Sobre la renovación de las cámaras, en el aula se señaló que el Senado de Estados Unidos se renueva por tercios y que la Cámara de Diputados se renueva completa, mientras que en la Argentina las cámaras se renuevan parcialmente.

\begin{note}
\textbf{Nota editorial de los apuntes originales.} Precisión: en la Argentina, la Cámara de Diputados se renueva por mitad cada dos años (art.~50, CN) y el Senado por tercios de los distritos cada dos años (art.~56, CN). En Estados Unidos, el Senado se renueva por tercios cada dos años y la Cámara de Representantes se renueva completa cada dos años.
\end{note}

\begin{formula}{Art.~50 --- Constitución Nacional}
«Los diputados durarán en su representación cuatro años, y son reelegibles; pero la Sala se renovará por mitad cada bienio; a cuyo efecto los nombrados para la primera Legislatura, luego que se reúnan, sortearán los que deban salir en el primer período.»
\end{formula}

\begin{formula}{Art.~56 --- Constitución Nacional}
«Los senadores duran seis años en el ejercicio de su mandato, y son reelegibles indefinidamente; pero el Senado se renovará a razón de una tercera parte de los distritos electorales cada dos años.»
\end{formula}

\begin{table}[H]
\centering
\begin{tabularx}{\textwidth}{>{\raggedright\arraybackslash}p{0.22\textwidth} >{\raggedright\arraybackslash}X >{\raggedright\arraybackslash}X}
\toprule
\textbf{Aspecto} & \textbf{Argentina} & \textbf{Estados Unidos} \\
\midrule
\textbf{Mandato de los diputados} & 4 años & 2 años \\
\textbf{Mandato de los senadores} & 6 años & 6 años (equivale a tres mandatos de diputado) \\
\textbf{Sistema electoral} & Representación proporcional (sistema D’Hondt) & Estados con sistema uninominal; menor proporcionalidad, bipartidismo y mayorías más afianzadas \\
\textbf{Perpetuación en el cargo} & La alternancia se exige para los cargos ejecutivos, no para los legislativos & Se discute la perpetuación de los senadores (20 o 30 años) \\
\bottomrule
\end{tabularx}
\end{table}

\subsection{Cierre del tema}

El fallo Formosa es valioso y debería brindar un principio orientador para todos los niveles de gobierno, en tanto fue establecido como estándar para las constituciones provinciales.

%----------------------------------------------------------
% CAPÍTULO 4
%----------------------------------------------------------
\chapter[Prácticas cuestionadas]{Prácticas cuestionadas}

\section{La ley de lemas}

\subsection{Planteo y antecedentes}

Santa Cruz tiene un punto recientemente modificado que existe en algunos estamentos: a determinado nivel de gobierno, la \textbf{ley de lemas}, otro cuestionamiento importante. La Corte se pronunció sobre el caso de Santa Fe, que también tenía ley de lemas, con un fallo más bien acotado, en el que apeló a una suerte de \textbf{zona de reserva provincial} para regular su propio sistema electoral.

Se impugnó la ley de lemas como violatoria del principio representativo, por tratarse de una norma que podía inducir a error y confusión al votante, pero la impugnación fue desestimada por la Corte. El docente considera que una nueva edición de este caso en la Corte Suprema podría tener un resultado distinto.

\subsection{Cómo funciona la ley de lemas}

La ley de lemas prevé distintas fórmulas, denominadas \textbf{sublemas}, dentro de una misma fuerza política. Todas las fórmulas de ese mismo lema aglutinan su número de votos, y resulta elegida, o representa ese espacio, la fórmula más votada dentro de los distintos sublemas. Se trata de un mecanismo que, cuanto menos, ofrece cierta opacidad y puede generar confusión en el electorado.

\begin{important}
\textbf{Ley de lemas.} Los votos de todos los sublemas se suman al lema y gana la fórmula más votada dentro de él. La Corte rechazó la impugnación en el caso de Santa Fe apelando a la zona de reserva provincial; el docente cree que hoy el resultado podría ser distinto por su opacidad y su potencial de confundir al electorado.
\end{important}

\section{Las candidaturas testimoniales}

Como antecedente parecido, las candidaturas testimoniales fueron cuestionadas y, previo a la elección, la Cámara Electoral las definió como un \textbf{mecanismo electoral ilegítimo}. Se trataba de candidatos políticos que figuraban a la cabeza de listas pero que tenían una incompatibilidad y de antemano se aceptaba que no ocuparían el cargo.

Pueden darse casos en los que el funcionario o congresista electo termina ejerciendo funciones en el Ejecutivo y abandona su candidatura antes de asumir. Aquí, en cambio, la situación se generaba a priori, con un rasgo estructural, y además existían figuras políticas muy prominentes que sí o sí iban a poder generar confusión y engañar al votante como cabezas de lista, sabiendo que tenían impedimentos, como los que tienen los gobernadores para integrar listas de legisladores. Eso está prohibido en la Constitución y alimentaba la candidatura testimonial. Por lo tanto, no solo confunde al votante, sino que puede confundirlo al punto de generar la ilusión de una candidatura expresamente prohibida en la Constitución, como la de un gobernador que sea legislador nacional.

\begin{formula}{Art.~73 --- Constitución Nacional}
«Los eclesiásticos regulares no pueden ser miembros del Congreso, ni los gobernadores de provincia por la de su mando.»
\end{formula}

\begin{important}
\textbf{Candidaturas testimoniales.} Postular como cabeza de lista a quien ya se sabe que no asumirá, o que tiene un impedimento constitucional, engaña al votante; la Cámara Electoral las consideró un mecanismo ilegítimo.
\end{important}

% TODO: contenido incompleto o ambiguo en las notas originales (la transcripción finaliza de forma abrupta; no se dispone del cierre ni de las eventuales consultas posteriores).

%----------------------------------------------------------
% CORNELL
%----------------------------------------------------------
\chapter[Cuadros Cornell de repaso]{Cuadros Cornell de repaso}

\section{Diseño constitucional del sistema político}

\begin{longtable}{@{}>{\raggedright\arraybackslash}p{0.27\textwidth}!{\vrule width 0.4pt}>{\raggedright\arraybackslash}p{0.67\textwidth}@{}}
\toprule
\textbf{Preguntas / palabras clave} & \textbf{Notas / respuestas} \\
\midrule
\endfirsthead
\toprule
\textbf{Preguntas / palabras clave} & \textbf{Notas / respuestas} \\
\midrule
\endhead

\textbf{¿Cuál es el diseño básico del sistema político y electoral?} &
Republicano, representativo y federal. Elección directa del Poder Ejecutivo y del Congreso. Elección indirecta de magistrados mediante el Consejo de la Magistratura; los jueces nacionales de la Capital aún se designan con ternas y acuerdo del Senado, aunque la justicia nacional se transfiere paulatinamente a la justicia local. \\
\cmidrule{1-2}

\textbf{¿Por qué la materia electoral es reserva del Congreso?} &
Para evitar que el poder de turno la modifique a su conveniencia. Requiere ley: el art.~99, inc.~3, CN veda los decretos de necesidad y urgencia en materia penal, tributaria, electoral y de partidos. Ejemplo: la ley de primarias no puede reformarse por decreto. \\
\midrule

\multicolumn{2}{@{}p{\dimexpr0.94\textwidth+2\tabcolsep\relax}@{}}{%
\textbf{Resumen:} El sistema político argentino es republicano, representativo y federal, y por eso la Constitución reserva la materia electoral y de partidos al Congreso, vedando los decretos de necesidad y urgencia.
} \\
\bottomrule
\end{longtable}

\section{Partidos políticos y sufragio}

\begin{longtable}{@{}>{\raggedright\arraybackslash}p{0.27\textwidth}!{\vrule width 0.4pt}>{\raggedright\arraybackslash}p{0.67\textwidth}@{}}
\toprule
\textbf{Preguntas / palabras clave} & \textbf{Notas / respuestas} \\
\midrule
\endfirsthead
\toprule
\textbf{Preguntas / palabras clave} & \textbf{Notas / respuestas} \\
\midrule
\endhead

\textbf{¿Qué rol cumplen los partidos políticos y qué exige la ley orgánica?} &
Son instituciones fundamentales del sistema democrático (art.~38 CN) y tienen la exclusividad de nominar candidatos. Deben tener personería jurídica, estatuto, afiliados mínimos por jurisdicción, y renovación y alternancia de autoridades. Con reconocimiento en cinco o más distritos, son partidos nacionales. Los controla la justicia electoral. \\
\cmidrule{1-2}

\textbf{¿Qué dijo el fallo Ríos?} &
Los partidos son el único vehículo institucional de nominación; no existe la candidatura independiente. Si cualquiera pudiera candidatearse por sí solo, el sistema sería fragmentado, atomizado e inoperante para representar al pueblo. \\
\cmidrule{1-2}

\textbf{¿Cada cuánto hay elecciones y por qué?} &
Cada dos años, por regla general, para asegurar la alternancia y el ejercicio continuo del sufragio, fuente de legitimidad de los poderes del Estado. \\
\cmidrule{1-2}

\textbf{¿Qué estableció Mignone (2002)?} &
Los procesados sin condena firme deben tener asegurado el voto; privarlos de él viola el principio de inocencia. Solo la condena efectiva, firme, por delito doloso con pena privativa de la libertad conlleva la inhabilitación. \\
\cmidrule{1-2}

\textbf{¿Cómo es el sufragio en Argentina?} &
Universal, secreto y obligatorio (art.~37 CN, que además lo califica de igual). Es fuente de legitimación de los cargos electivos y su obligatoriedad no genera disensos entre los partidos. \\
\midrule

\multicolumn{2}{@{}p{\dimexpr0.94\textwidth+2\tabcolsep\relax}@{}}{%
\textbf{Resumen:} Los partidos políticos, controlados por la justicia electoral, son el único vehículo de nominación de candidatos (Ríos), y el sufragio universal, secreto y obligatorio es la fuente de legitimidad del poder; ese derecho alcanza incluso a quienes están detenidos sin condena firme (Mignone).
} \\
\bottomrule
\end{longtable}

\section{Alternancia y reelección}

\begin{longtable}{@{}>{\raggedright\arraybackslash}p{0.27\textwidth}!{\vrule width 0.4pt}>{\raggedright\arraybackslash}p{0.67\textwidth}@{}}
\toprule
\textbf{Preguntas / palabras clave} & \textbf{Notas / respuestas} \\
\midrule
\endfirsthead
\toprule
\textbf{Preguntas / palabras clave} & \textbf{Notas / respuestas} \\
\midrule
\endhead

\textbf{¿Qué es la alternancia y por qué importa?} &
Es la renovación de quienes ejercen el poder. La jurisprudencia la destaca como signo de afianzamiento de los sistemas democráticos, representativos y republicanos y como base de los frenos y contrapesos. \\
\cmidrule{1-2}

\textbf{¿Qué resolvió la Corte en el caso Formosa?} &
Declaró inconstitucional el art.~132 de la Constitución provincial (reelección indefinida), apoyándose en el sistema republicano y en la OC-28/21 de la Corte IDH. No fijó un número exacto de períodos, pero admitió que hasta dos reelecciones podrían tolerarse; el período de más de 30 años socavaba la alternancia mediante una «superestructura política permanente». \\
\cmidrule{1-2}

\textbf{¿Se aplica la doctrina Formosa a intendentes y otras provincias?} &
Es probable, pero haría falta un nuevo fallo. Antecedentes: Zamora (Santiago del Estero, 2013) y Manzur (Tucumán, 2023). Catamarca y Santa Cruz no tienen límite; San Juan prevé tres mandatos. El principio de alternancia debería regir en todo cargo ejecutivo. \\
\cmidrule{1-2}

\textbf{¿Rige la alternancia para los legisladores? ¿Y en EE.\,UU.?} &
No para los cargos legislativos. En EE.\,UU. se discute la perpetuación de los senadores (20 o 30 años); allí los representantes duran 2 años y los senadores 6, frente a 4 y 6 en Argentina. \\
\midrule

\multicolumn{2}{@{}p{\dimexpr0.94\textwidth+2\tabcolsep\relax}@{}}{%
\textbf{Resumen:} La alternancia es el hilo conductor: el fallo Formosa muestra que una reelección formalmente legítima pero indefinida socava la república, y abre el debate sobre intendentes y otras provincias.
} \\
\bottomrule
\end{longtable}

\section{Prácticas cuestionadas}

\begin{longtable}{@{}>{\raggedright\arraybackslash}p{0.27\textwidth}!{\vrule width 0.4pt}>{\raggedright\arraybackslash}p{0.67\textwidth}@{}}
\toprule
\textbf{Preguntas / palabras clave} & \textbf{Notas / respuestas} \\
\midrule
\endfirsthead
\toprule
\textbf{Preguntas / palabras clave} & \textbf{Notas / respuestas} \\
\midrule
\endhead

\textbf{¿Qué es la ley de lemas?} &
Distintas fórmulas (sublemas) dentro de una misma fuerza suman sus votos y gana la fórmula más votada del lema. Es opaca y puede confundir al votante. La Corte la convalidó en Santa Fe por la zona de reserva provincial, pero hoy podría decidir distinto. \\
\cmidrule{1-2}

\textbf{¿Qué son las candidaturas testimoniales?} &
Candidatos a la cabeza de listas que se sabe de antemano que no asumirán o tienen un impedimento (por ejemplo, gobernadores como legisladores nacionales). La Cámara Electoral las consideró un mecanismo ilegítimo porque engañan al votante. \\
\midrule

\multicolumn{2}{@{}p{\dimexpr0.94\textwidth+2\tabcolsep\relax}@{}}{%
\textbf{Resumen:} La ley de lemas y las candidaturas testimoniales ilustran mecanismos que, por su opacidad, pueden confundir al votante y deteriorar la representación política.
} \\
\bottomrule
\end{longtable}

\end{document}
